\renewcommand{\titreDiapos}{Essai}% fallback definition

\section{Sujet du cours}


\begin{frame}{Le sujet, en bref}

\begin{defblock}{Documents, collections massives, stockage et calcul distribués}
Représentation d'entités sous forme de  \alert{documents}
pour la gestion de \alert{collections} à très grande échelle dans des environnements distribués.
\end{defblock}

\vfill

Les thèmes traités\,:

\begin{redbullet}
  \item \alert{modélisation} sous forme  de \alert{document}, structuré, semi-structuré, non structuré
  \item \alert{recherche exacte} et  \alert{recherche approchée} dans de très grandes collections
  \item \alert{traitements par lot} (MapReduce et +)
  \item \alert{distribution du stockage} (réplication, partitionnement).
  \item \alert{distribution des calculs}
\end{redbullet}

\vfill
\begin{blocgauche}
Des mots-clés à comprendre (et à critiquer): données structurées/non structurées, bases NoSQL, Cloud, BigData,
moteurs de recherche.
\end{blocgauche}
\end{frame}


\begin{frame}
\frametitle{Point de départ\,: bases relationnelles}

Les \alert{bases relationnelles} = applications
gérant des informations \alert{structurées} et \alert{régulières}\,: 
applications de ``gestion'',
 Web,  mobiles.

\vfill

\begin{redbullet}
  \item Une modélisation normalisée
  \item Un langage (SQL) très bien défini, normalisé lui aussi
  \item De très bonnes performances, obtenues \alert{automatiquement}
  \item Une gestion robuste de la concurrence d'accès
\end{redbullet}


\vfill

\alert{Attention à bien apprécier ce qu'on gagne / perd en passant au "NoSQL"} (on gagne
marginalement, on perd beaucoup)

\vfill
\begin{blocgauche}
Révisions nécessaires? Voir  \url{http://sql.bdpedia.fr} et \url{http://sys.bdpedia.fr}
\end{blocgauche}

\end{frame}

\begin{comment}
\begin{frame}
\frametitle{Problématique 1\,: La notion de ``document''}

Document = unité d'information autonome ou quasi-autonome.

\begin{redbullet}
  \item Peu ou pas de référence à d'autres documents.
  
\item Peu ou pas de structure; ou une structure très flexible.

\item Un contenu souvent à orientation multimédia.

\end{redbullet}

\alert{Examples (1)}: documents textuels, types documents Web.
   
\vfill
   
\alert{Examples (2)}:  images,  documents
audios,  vidéos\,; pas de structure explicite, production de descripteurs synthétiques pour
   tenter de les indexer.

\vfill

\alert{Examples (3)}: jeux en ligne\,: artifacts graphiques, actions
utilisateur.


\vfill

\alert{Examples (4)}: tous les fichiers de votre ordinateur...

\vfill

\begin{blocgauche}
Impose de repenser la notion de \alert{modèle}, schéma et représentation des données.
\end{blocgauche}
\end{frame}


\begin{frame}
\frametitle{Problématique 2\,: recherche d'information}

Dans les collections documentaires, peu ou pas de structure fixe,

\begin{redbullet}
  \item SQL inadapté.
  \item Recherche \guill{exacte} souvent insatisfaisante.
\end{redbullet}

\medskip

La recherche
s'effectue souvent \alert{par similarité}

\begin{redbullet}
  \item on fournit un document ``requête''.
  \item le système recherche les documents \alert{proches} du document-requête.
 \end{redbullet}

\vfill

\begin{blocgauche}
Par exemple, quand on recherche sur le Web:


\begin{redbullet}
  \item Un ensemble de mots-clés: c'est le document requête\,;
  \item le moteur de recherche trouve les documents les plus proches (on verra comment)\,;
 \end{redbullet}
\end{blocgauche}

\end{frame}

\begin{frame}
\frametitle{Problématique 3\,: stockage à très grande échelle}

On atteint facilement des \alert{volumes} extrêmement importants (Téraoctets+)

\begin{redbullet}
  \item les moteurs de recherche qui collectent des \alert{documents}
  disponibles sur le Web.
  \item  les applications utilisées à l'échelle du Web\,; commerce
  électronique (Amazon); réseaux sociaux (Facebook).
  \item  données gérées par les jeux en ligne.
\end{redbullet}


\vfill

\begin{blocgauche}
 Nouveaux systèmes, dits ``NoSQL'' sont conçus pour gérer de vastes collections
 de documents de manière scalable
 \begin{redbullet}
    \item pour les accès temps réel;
    \item pour les traitements analytiques (MapReduce).
 \end{redbullet}
\end{blocgauche}

\end{frame}


\begin{frame}
\frametitle{Problématique 4\,: calcul  à très grande échelle}

Des environnements de développement et d'exécution
\alert{distribués} pour traiter les grands volumes de données.


 \begin{redbullet}
    \item traitement de flux;
    \item traitement de grands masses de données statiques.
 \end{redbullet}
 
\vfill


À chaque fois, les principes, et une illustration  avec les systèmes.
\begin{redbullet}
  \item MapReduce
  \item  Traitement de flux avec FLink
  \item  Traitement de données massives avec Spark
\end{redbullet}


\begin{blocgauche}
\end{blocgauche}

\end{frame}
\end{comment}
%%%%%%%%%%%%%%%%%%%%%
\section{Objectifs}


\begin{frame}
\frametitle{Plan du cours\,: deux parties, $\approx$ 6 semaines chacune}

\alert{Partie 1}\,: modèles de représentation et de traitement

\begin{redbullet}
  \item \alert{Représentation de documents textuels}\,: JSON; la modélisation
    de documents, les langages 
     de manipulation

  \item \alert{Recherche exacte et approchée}\,:  langages d'interogation,
      principes de la recherche d'information

  \item \alert{Les traitements par lot de données massives}\,:  MapReduce et autres
          opérateurs de traitement distribués
 \end{redbullet}

\alert{Partie 2}\,: systèmes distribués
 
\begin{redbullet}    
   \item \alert{Stockage distribué}. Partition, réplication, reprise sur panne
   
   \item \alert{Calcul distribué}. Les chaînes de traitement distribué
\end{redbullet}

\vfill
Mise en pratique: des TPs avec Cassandra, ElasticSearch et Spark

\end{frame}


\section{Prérequis}

\begin{frame}{Prérequis}

\begin{redbullet}
   \item \alert{Compréhension des bases relationnelles}, soit au moins la conception d'un schéma, 
      SQL, ce qu'est un index et des notions de base sur les transactions.
      
      $\Rightarrow$ si nécessaire, voir \url{http://www.bdpedia.fr}
      
   \item \alert{Une aisance minimale dans un environnement de développement}. Editer un fichier, lancer une commande, 
      ne pas paniquer devant un nouvel outil, savoir résoudre un problème avec un minimum de tenacité, etc.      

    $\Rightarrow$   il est utile (mais pas \alert{obligatoire}) reproduire les exemples donnés.
      
\end{redbullet}

\begin{blocgauche}
Il ne s'agit pas \alert{d'apprendre} des systèmes, mais de comprendre 
des principes via la pratique (limitée)
\end{blocgauche}

\end{frame}


\subsection{Environnement de travail}
\begin{frame}{Environnement matériel et logiciel}

\alert{Pour la mise en pratique}\,: 

\begin{redbullet}
  \item au Cnam, tout vous est fourni;
  \item Cnam à distance, ou en auditeur libre, il vous faut une machine  dotée d'au moins 8 GO de mémoire (Linux, Mac/OS ou Windows).
\end{redbullet}

\vfill

\alert{Tous les logiciels utilisés sont libres de droits}. Nous travaillons essentiellement avec
\begin{redbullet}
   \item  \href{http://elasticsearch.org}{ElasticSearch} 
      pour l'indexation de documents.
   \item \href{http://cassandra.apache.org}{Cassandra} 
   \item \href{http://spark.apache.org}{Spark}
   \end{redbullet}

\vfill

\begin{blocgauche}
L'installation de ces outils (et leur utilisation) peut se faire avec \alert{Docker}, un
gestionnaire de systèmes distribués virtuels.
\end{blocgauche}
\end{frame}


\section{Organisation}
\begin{frame}{Comment travailler}

Le cours est découpé en \alert{chapitres}, couvrant un sujet bien déterminé, et en \alert{sessions}.

\smallskip

\alert{L'unité de travail est la session}, correspondant à 1-2 heures de travail personnel

\vfill


Pour assimiler une session vous pouvez combiner les ressources suivantes:

\begin{redbullet}
  \item \alert{La lecture} du \href{http://b3d.bdpedia.fr}{support en ligne}, 
       également disponible en \href{http://b3d.bdpedia.fr/files/b3d.pdf}{PDF} ou en 
       \href{http://b3d.bdpedia.fr/files/b3d.epub}{ePub}.
    \item \alert{Le suivi du cours}, en vidéo ou en présentiel.
   \item \alert{Le test des exemples de code} fournis dans chaque session.
   \item \alert{La réalisation des exercices} proposés en fin de session.
\end{redbullet}

\vfill
\begin{blocgauche}
Pendant les visios\,: reprise (rapide) des principes, et nous
faisons les quiz.
\end{blocgauche}

\end{frame}

\begin{frame}{Mener votre travail personnel}

Vous devez maîtriser le contenu des sessions \alert{dans l'ordre
où elles sont proposées}.

\vfill

\alert{Commencez par lire le support}, jusqu'à ce que les principes vous paraissent clairs.

\vfill

\alert{Suivez la visio}, pour une synthèse et un média complémentaire

\vfill

Reproduisez les exemples de code, \alert{sans y passer trop de temps}

\vfill

\begin{blocgauche}
\alert{Cherchez
à résoudre les problèmes par vous-mêmes} au besoin: c'est le meilleur moyen de comprendre.

\vfill

\alert{Finissez par les exercices}. Faites les vous-mêmes, avant de regarder la solution.

\vfill

\alert{Préparez les quiz} proposés en fin des principaux chapitres.
\end{blocgauche}

\end{frame}

